% Source: https://artofproblemsolving.com/community/c3h2558989_difference_of_5th_power
% Stable identity: aops-topic-2558989-q01
Setting $x=y$ makes $x^5-y^5$ zero, so $x-y$ is a factor. Imitating the sum of a geometric series, try
\[(x-y)\left(x^4+x^3y+x^2y^2+xy^3+y^4\right).\]
Expanding,
\begin{align*}
&x\left(x^4+x^3y+x^2y^2+xy^3+y^4\right)-y\left(x^4+x^3y+x^2y^2+xy^3+y^4\right)\\
&=\left(x^5+x^4y+x^3y^2+x^2y^3+xy^4\right)-\left(x^4y+x^3y^2+x^2y^3+xy^4+y^5\right)=x^5-y^5,
\end{align*}
since all the middle terms cancel in pairs. Hence
\[\boxed{x^5-y^5=(x-y)\left(x^4+x^3y+x^2y^2+xy^3+y^4\right)}.\]
For example, $x=2$, $y=1$ gives $31=1\cdot(16+8+4+2+1)$.
